\section{ÁREAS DE ATUAÇÃO}

O \NOMEEGRESSO é um profissional de nível superior capacitado a planejar, projetar, implementar, administrar e gerenciar sistemas de automação industrial, considerando rigor técnico, fundamentos científicos e responsabilidade socioambiental. Atua com postura empreendedora e consciência de seu papel no contexto político, econômico e produtivo.

Trata-se de um profissional com formação generalista, apto a intervir no controle e na automação de equipamentos, processos, unidades e sistemas de produção. Sua atuação caracteriza-se pela interface direta entre os sistemas produtivos e os sistemas gerenciais das organizações, planejando e executando soluções de controle e automação orientadas ao aumento da eficiência, confiabilidade e segurança industrial.

O egresso possui competência para estudar, projetar e especificar materiais, componentes, dispositivos e equipamentos elétricos, eletromecânicos, eletrônicos, ópticos, magnéticos e de instrumentação; bem como para planejar, instalar, operar e manter sistemas de medição, acionamento, controle e supervisão de processos industriais. Também atua no projeto, instalação e manutenção de robôs industriais, sistemas de manufatura e redes industriais, podendo coordenar equipes, conduzir estudos de viabilidade técnico-econômica, supervisionar serviços técnicos e elaborar laudos, pareceres e avaliações. Em todas essas atividades, observa rigor ético, atendimento às normas, prudência técnica, eficiência energética e mitigação de impactos ambientais.

O egresso do \NOMECURSO está habilitado a atuar em empresas de engenharia, indústrias fabricantes de equipamentos e software para automação, bem como nos diversos setores usuários de sistemas automáticos. Entre suas áreas de intervenção destacam-se:

\begin{enumerate}
	\item Automação de processos e sistemas nos setores industrial, comercial, residencial e de serviços;
	\item Modernização, otimização e manutenção de unidades de produção automatizadas;
	\item Projeto e integração de sistemas de automação industrial;
	\item Concepção, instalação e operação de unidades de produção automatizadas;
	\item Desenvolvimento de soluções para instrumentação, controle, supervisão e gestão de processos industriais;
	\item Treinamento e capacitação de equipes técnicas em indústrias e instituições de ensino;
	\item Pesquisa e desenvolvimento de tecnologias aplicadas ao controle e automação.
\end{enumerate}

Nesse contexto, o Engenheiro de Controle e Automação encontra oportunidades em concessionárias de energia, indústrias de base e transformação, plantas de manufatura automatizada, automação predial, setores de petróleo e gás, siderurgia, têxtil, calçadista e parques industriais diversos. Sua presença é essencial tanto na modernização de processos produtivos consolidados quanto na implantação de soluções tecnológicas alinhadas às demandas contemporâneas de produtividade, segurança operacional e sustentabilidade.

